\begin{frame}[shrink=12]\frametitle{Результаты сравнительных измерений}
{\footnotesize\textit{ART: 16~шардов, без сортировки при выгрузке · Ryzen~AI~9~HX~370, 24~потока, 32~ГБ LPDDR5X}}

\vspace{0.3em}
\setlength{\tabcolsep}{3pt}\tiny
\resizebox{\linewidth}{!}{%
\begin{tabular}{p{0.75cm}|p{3.2cm}|p{2.2cm}|p{2.2cm}}
\hline
\rowcolor{gray!20}
\textbf{Набор} & \textbf{Операция} & \textbf{ranger} & \textbf{шард. ART} \\
\hline\hline
$10^5$ & Полная выгрузка IPv4 & 0.24--0.28~мс & 11--12~мс \\\hline
$10^5$ & Полная выгрузка IPv6 & 0.20--0.25~мс & 11--12~мс \\\hline
$10^5$ & Одиночная запись     & 590--620~мс  & ${\approx}1{,}7{\cdot}10^{-4}$~мс \\\hline
$10^5$ & Поиск IPv4 /32       & 151~нс       & 400--500~нс \\\hline
$10^5$ & Поиск IPv6 /128      & 168--178~нс  & 3700--4000~нс \\\hline
\rowcolor{gray!12}
3M & Полная загрузка таблицы & 27{,}423~с & 0{,}622~с \\\hline
\rowcolor{gray!12}
3M & Полная выгрузка снимка & 7{,}522~мс & 326{,}957~мс \\\hline
\rowcolor{gray!12}
3M & Поиск по IP & 185{,}3~нс & 2{,}899~мкс \\\hline
\rowcolor{gray!12}
3M & Одиночная запись & 28{,}178~с & 156{,}8~нс \\\hline
\end{tabular}
}

\vspace{0.3em}
\footnotesize
\textbf{Вывод:} ART выигрывает на полной загрузке и одиночной записи; \texttt{ranger} быстрее на полной выгрузке и поиске. Для real-time ключевым остаётся локальное применение изменений без пересборки всей таблицы.
\end{frame}
